\section{Kathrin's stuff}

\khnote{In this section I recall some background from previous literature:
	I recall definitions and notions for PKE schemes and KEMs and the Fujisaki--Okamoto transformation.
}

\subsection{Public-Key Encryption Schemes}\label{app:FO-etc:PKE}
	
	\begin{definition}[Public-Key Encryption Scheme]\label{def:PKE}
		A public-key encryption scheme $\PKE = (\KG, \Enc, \Dec)$ is a triple of algorithms
		defined with respect to public-key space \PKSpace, secret-key space \SKSpace,
		message (plaintext) space \MSpace, ciphertext space \CSpace,
		and randomness space \RSpace.
		More precisely, the triple of algorithms is defined as follows.
		\begin{itemize}
			\item
				$\KG() : (\pk, \sk) \in \PKSpace \times \SKSpace$.
				Probabilistic algorithm that takes no input and outputs a ``key pair'' $(\pk, \sk)$.
			\item
				$\Enc(\pk \in \PKSpace, m \in \MSpace) : c \in \CSpace$.
				Probabilistic algorithm that, given a public key $\pk$ and message $m$,
				produces an ``encryption'' (``ciphertext'') $c$ of $m$ under $\pk$.
				We may explicitly specify the randomness used in encryption by writing $\Enc(\pk, m; r)$,
				where $r \in \RSpace$.
			\item
				$\Dec(c \in \CSpace, \sk \in \SKSpace) : m \in \MSpace \cup \{\bot\}$.
				Deterministic algorithm that, given a secret key $\sk$ and ciphertext $c$,
				yields either a message $m \in \MSpace$ or an explicit failure symbol $\bot \notin \MSpace$,
				the latter indicating that $c$ is not a valid ciphertext (under $\sk$).
		\end{itemize}
	\end{definition}
	
	\todo{OWCPA or INDPCA, correctness, rigidity}

\subsection{Key Encapsulation Mechanisms}\label{app:FO-etc:KEMs}
	
	\begin{definition}[Key Encapsulation Mechanism]\label{def:KEM}
		A key encapsulation mechanism $\KEM = (\KG, \Encaps, \Decaps)$ is a triple of algorithms
		defined with respect to public-key space \PKSpace, secret-key space \SKSpace,
		key space \KSpace, and ciphertext space \CSpace.
		More precisely, the triple of algorithms is defined as follows.
		\begin{itemize}
			\item
				$\KG() : (\pk, \sk) \in \PKSpace \times \SKSpace$.
				Probabilistic algorithm that takes no input and outputs a ``key pair'' $(\pk, \sk)$.
			\item
				$\Encaps(\pk \in \PKSpace) : (K, c) \in \KSpace \times \CSpace$.
				Probabilistic algorithm that, given a public key $\pk$,
				produces a tuple $(c,\KEMoutputKey)$,
				where $c$ is said to be an ``encapsulation'' (``ciphertext'') of the key $\KEMoutputKey$.
			\item
				$\Decaps(\sk \in \SKSpace, c \in \CSpace) : \KEMoutputKey \in \KSpace \cup \{\bot\}$.
				Deterministic algorithm that, given a secret key $\sk$ and ciphertext $c$,
				yields either a key $\KEMoutputKey \in \KSpace$ or an explicit failure symbol $\bot \notin \KSpace$,
				the latter indicating that $c$ is not a valid ciphertext (under $\sk$).
		\end{itemize}
	\end{definition}
	
	We also recall the standard security definition for KEMs,
	\underline{Ind}istinguishability under \underline{C}hosen \underline{C}iphertext \underline{A}ttacks ($\INDCCA$).
	
	\begin{definition}[\INDCCA]\label{def:indcca}
		We define the \INDCCA game as in Figure~\ref{fig:def:IND-CCA}
		and the \INDCCA advantage function of an adversary \AdversaryA (with binary output) against KEM as 
		\[
			\Advantage_{\KEM}^{\INDCCA}(\AdversaryA)
			= \left| \Pr\left[\INDCCA^{\AdversaryA}_\KEM \Rightarrow 1\right] - \frac{1}{2} \right| \enspace .
		\]
		
		
		\begin{figure}[ht]
			\begin{center}
				\fbox{
					\small
					\nicoresetlinenr
					\begin{minipage}[t]{5cm}	
						\underline{{\bf Game} $\INDCCA_{\KEM}$:}
						\begin{nicodemus}
							\item $(\pk, \sk) \leftarrow \KG$
							\item $b \sample \{0,1\}$
							\item $(K_0,c^*) \leftarrow \Encaps(\pk)$
							\item $K_1 \sample \KSpace$
							\item $b' \leftarrow A^{\oracleDecaps}(\pk, c^*, K_b)$
							\item \pcreturn $\bool{b' = b^*}$
						\end{nicodemus}
					\end{minipage}
					
					\quad
					
					\begin{minipage}[t]{3.5cm}
						\underline{ $\oracleDecaps( c\ne c^*)$:}
						\begin{nicodemus}
							\item $K \gets \Decaps(\sk,c)$
							\item \pcreturn $K$
						\end{nicodemus}	
					\end{minipage}
				}
			\end{center}
			\caption{\INDCCA game for $\KEM = (\KG, \Encaps, \Decaps)$.}
			\label{fig:def:IND-CCA}
		\end{figure}
		
	\end{definition}
						
				
\subsection{The Fujisaki-Okamoto Transform}

	In this section, we recall the definition of one of the \FO-transforms that were given in~\cite{TCC:HofHovKil17},
	as the composition of the two following transformations:
	
	\begin{itemize}
		\item
			The de-randomizing \emph{\T-transform} that additionally adds a re-encryption check to the decryption procedure;
			and 
		\item
			PKE-to-KEM \emph{\U-transform}\todo{Specify variant} that derives a KEM output key \KEMoutputKey
			from a randomly chosen message $m$, which it encrypts using \PKE.
			(\cite{TCC:HofHovKil17} also defined other \U-variants that differ from the variant discussed here
			in how they handle invalid ciphertexts and in what they hash into the output key \KEMoutputKey.)
	\end{itemize} 
	
	

	\begin{definition}[\T-transform]
		To a PKE scheme $\PKE = (\KG, \Enc, \Dec)$ and a hash function $\ROrand:\MSpace \rightarrow \RSpace$,
		we associate the (deterministic) PKE scheme $\T[\PKE, \ROrand]$ defined in Figure~\ref{fig:def:T-transform}.
		
		\begin{figure}[t]
			\begin{center}
				\fbox{
					\small
					\nicoresetlinenr
					
					\begin{minipage}[t]{3.6cm}	
						\underline{ $\TEnc(pk, m)$:}
						\begin{nicodemus}
							\item $c \gets \Enc(pk, m; \ROrand(m))$
							\item \pcreturn $c$
						\end{nicodemus}
					\end{minipage}
					
					\quad
					
					\begin{minipage}[t]{5cm}
						\underline{ $\TDec(sk, c)$:}
						\begin{nicodemus}
							\item $m^\prime \gets \Enc(sk,c)$
							\item \pcif $m^\prime = \perp$ or $c \neq \Enc(pk, m; \ROrand(m))$
							\item \quad \pcreturn $\perp$
							\item \pcelse \pcreturn $m^\prime$
						\end{nicodemus}	
					\end{minipage}
				}
			\end{center}
			\caption{Algorithms of $\T[\PKE, \ROrand]=: (\KG, \TEnc, \TDec)$.}
			\label{fig:def:T-transform}
		\end{figure}
		
	\end{definition}

	
	
	\begin{definition}[\UImpMess-transform]
		%
		To a PKE scheme $\PKE = (\KG, \Enc, \Dec)$ and a hash function \ROkey,
		we associate the KEM $\KEMImpMess := \UImpMess[\PKE, \ROkey]$ defined in~\ref{fig:def:U-transform}.
		
		\begin{figure}[ht]
			\begin{center}
				\fbox{
					\small
					\nicoresetlinenr
					\begin{minipage}[t]{3cm}
						\underline{$\KGImp()$}
						\begin{nicodemus}
							\item $(pk, sk) \sample \KG()$
							\item $\rejectionSeed \sample \MSpace$
							\item $\sk' \leftarrow (sk, \rejectionSeed)$
							\item \pcreturn $(pk, \sk')$
						\end{nicodemus}
					\end{minipage}
					
					\quad
					
					\begin{minipage}[t]{3.2cm}
						\underline{$\EncapsMess(pk)$}
						\begin{nicodemus}
							\item $m \sample \MSpace$
							\item $c \leftarrow \Enc(pk, m)$
							\item $\KEMoutputKey := \ROkey(m)$
							\item \pcreturn $(\KEMoutputKey, c)$
						\end{nicodemus}
					\end{minipage}
					
					\quad
					
					\begin{minipage}[t]{3.4cm}
						\underline{$\DecapsImpMess(\sk', c)$}
						\begin{nicodemus}
							\item Parse $\sk' := (sk, \rejectionSeed)$
							\item $m' := \Dec(sk, c)$
							\item \pcif $m' = \bot$
								\item \pcreturn $\ROkey(\rejectionSeed, c)$
							\item \pcreturn $\ROkey(m')$
						\end{nicodemus}
					\end{minipage}
				}
			\end{center}
			\caption{Algorithms of $\UImpMess[\PKE, \ROkey, \ROimp] =: (\KGImp, \EncapsMess, \DecapsImpMess)$.}
			\label{fig:def:U-transform}			
		\end{figure} 
		
	\end{definition}